\documentclass[11pt,a4paper]{article}
\usepackage{fontspec,amsmath,amssymb,amsthm,microtype}
\setmainfont{lmroman10-regular.otf}[BoldFont=lmroman10-bold.otf,ItalicFont=lmroman10-italic.otf,BoldItalicFont=lmroman10-bolditalic.otf]
\usepackage{xcolor}
\usepackage[most]{tcolorbox}
\usepackage{mathrsfs}
\providecommand{\lt}{<}
\providecommand{\gt}{>}
\setmonofont{DejaVuSansMono.ttf}[BoldFont=DejaVuSansMono-Bold.ttf,ItalicFont=DejaVuSansMono-Oblique.ttf,BoldItalicFont=DejaVuSansMono-BoldOblique.ttf,Scale=0.86]
\newfontfamily\leanfallback{STIXTwoMath-Regular.otf}[Scale=0.86]
\definecolor{leanrule}{gray}{0.62}
\definecolor{leanback}{gray}{0.962}
\definecolor{leantitle}{gray}{0.22}
\definecolor{leanmuted}{gray}{0.38}
\newtcolorbox{leanbox}[1]{enhanced,breakable,colback=leanback,colframe=leanrule,
  boxrule=0.4pt,arc=0.6mm,left=2.4mm,right=2.4mm,top=1.8mm,bottom=1.8mm,
  before skip=0.8em,after skip=1em,fonttitle=\small\bfseries,coltitle=leantitle,
  colbacktitle=leanback,titlerule=0.3pt,title={#1}}
\newlength{\leancharwidth}
\newenvironment{leancode}{\par\smallskip\begingroup\ttfamily\footnotesize\raggedright
  \settowidth{\leancharwidth}{M}\setlength{\parskip}{0pt}\setlength{\parindent}{0pt}
  \frenchspacing}{\par\endgroup\smallskip}
\newcommand{\leanline}[2]{\hangindent=\dimexpr#1\leancharwidth+4\leancharwidth\relax
  \hangafter=1\noindent\hspace*{#1\leancharwidth}#2\par}
\newcommand{\leansource}[1]{\par{\footnotesize\color{leanmuted}#1}\par}
\newcommand{\leanlabel}[1]{\par\smallskip\noindent{\small\bfseries #1}\par}
\usepackage[margin=25mm]{geometry}
\usepackage[unicode,hidelinks]{hyperref}
\hypersetup{pdfauthor={LangevinLean Project}}
\setlength{\emergencystretch}{3em}
\setlength{\parskip}{0.4em}
\setlength{\parindent}{0pt}
\allowdisplaybreaks[2]
\raggedbottom
\title{\texorpdfstring{The conditional common-noise convergence theorem}{The conditional common-noise convergence theorem}}
\author{LangevinLean Project}
\date{}
\begin{document}\maketitle
Draft mathematical exposition; these notes may contain mistakes.


Let \(X_k\) be the numerical iterates and \(Y_k\) a supplied reference on the same probability space, with \(\mathcal L(Y_k)=\pi\). Assume a common-noise update contraction with factor \(0\le q<1\), a finite initial mean-square difference, and square-integrable reference defects bounded in norm by \(b\). Subtracting the updates gives \(e_{k+1}\le qe_k+b\), where \(e_k=\|X_k-Y_k\|_2\). Moment propagation ensures these norms are finite.



The joint law of \((X_n,Y_n)\) bounds \(W_2(\mathcal L(X_n),\pi)\) by \(e_n\). Iteration gives the stated transport estimate. Its initial term is the cost of the supplied coupling. Replacing that cost by an initial Wasserstein distance requires the same construction and constants for every finite-cost initial coupling.


\section{\texorpdfstring{Transport bound from a supplied discrete reference}{Transport bound from a supplied discrete reference}}


Let \(\nabla U\) be \(L\)-Lipschitz, \(0\le q<1\), and \(b\ge0\). Assume the initial vector, all noise vectors, and all reference vectors \(Y_k\) are measurable; \((Y_k)_\#P=\pi\); and

\[Y_{k+1}=T_h(U,h,Y_k,\xi_k)+V_k\quad\text{almost everywhere}.\]

Suppose the common-noise update is \(q\)-contractive, \(X_0-Y_0\in L^2\), every \(V_k\in L^2\), and \(\|V_k\|_2\le b\). Then

\[W_2((X_n)_\#P,\pi)\le q^n\|X_0-Y_0\|_2+\frac b{1-q}.\]

The initial term is the cost of the chosen coupled pair.



Let the iterates be \(X_k\) and let \(e_k=\|X_k-Y_k\|_2\).



Finite moments propagate from the initial error and the defects, so each \(e_k\) is a finite \(L^2\) norm. Every numerical iterate is measurable.



To derive the recurrence, fix \(k\) and set \(A=T_h(U,h,X_k,\xi_k)-T_h(U,h,Y_k,\xi_k)\).



Measurability of the gradient and the inputs makes \(A\) measurable, hence almost everywhere strongly measurable.



Subtract the two updates on the full-measure set where the reference identity holds. This yields \(X_{k+1}-Y_{k+1}=A-V_k\).



The common-noise contraction gives \(\|A\|\le q\|X_k-Y_k\|\).



The finite-moment synchronous bound and \(\|V_k\|_2\le b\) give the recurrence \(e_{k+1}\le qe_k+b\).



The joint law of the coupled pair bounds transport by \(e_n\). Its reference marginal is \(\pi\) by the supplied stationary law.



Iterate the derived recurrence to get \(e_n\le q^ne_0+b/(1-q)\), and compose with the transport comparison.



For the chosen reference, let \(e_k=(\mathbb E\|X_k-Y_k\|^2)^{1/2}\). The shared-noise decomposition gives \(e_{k+1}\le qe_k+b\), with all norms finite by the preceding induction. Therefore \(e_n\le q^ne_0+b(1-q^n)/(1-q)\). The coupling inequality and stationarity of \(Y_n\) yield the displayed result after replacing the finite geometric factor by its uniform upper bound.

\begin{leanbox}{Lean: \texttt{ula\_\allowbreak{}wasserstein\_\allowbreak{}bound\_\allowbreak{}of\_\allowbreak{}reference}}
\leansource{\texttt{ula\_\allowbreak{}wasserstein\_\allowbreak{}bound\_\allowbreak{}of\_\allowbreak{}reference} \quad UnadjustedLangevin/\allowbreak{}Convergence.lean, lines 25–72}
\begin{leancode}
\leanline{0}{theorem\ ula\_wasserstein\_bound\_of\_reference}
\leanline{4}{\{Ω\ :\ Type*\}\ [MeasurableSpace\ Ω]\ \{d\ :\ ℕ\}}
\leanline{4}{(P\ :\ Measure\ Ω)\ (U\ :\ E\ d\ →\ ℝ)\ (h\ :\ ℝ)\ (L\ q\ :\ ℝ≥0)\ (b\ :\ ℝ)}
\leanline{4}{(hLip\ :\ LipschitzWith\ L\ (gradient\ U))\ (hq1\ :\ (q\ :\ ℝ)\ <\ 1)\ (hb\ :\ 0\ ≤\ b)}
\leanline{4}{(initial\ :\ Ω\ →\ E\ d)\ (noise\ :\ ℕ\ →\ Ω\ →\ E\ d)}
\leanline{4}{(Y\ V\ :\ ℕ\ →\ Ω\ →\ E\ d)\ (π\ :\ Measure\ (E\ d))}
\leanline{4}{(hinitial\ :\ Measurable\ initial)\ (hnoise\ :\ ∀\ n,\ Measurable\ (noise\ n))}
\leanline{4}{(hY\ :\ ∀\ n,\ Measurable\ (Y\ n))\ (hlaw\ :\ ∀\ n,\ P.map\ (Y\ n)\ =\ π)}
\leanline{4}{(hreference\ :\ ∀\ n,\ ∀ᵐ\ ω\ ∂P,}
\leanline{6}{Y\ (n\ +\ 1)\ ω\ =\ ulaUpdate\ U\ h\ (Y\ n\ ω)\ (noise\ n\ ω)\ +\ V\ n\ ω)}
\leanline{4}{(hcontract\ :\ ∀\ x\ y\ z,}
\leanline{6}{‖ulaUpdate\ U\ h\ x\ z\ -\ ulaUpdate\ U\ h\ y\ z‖\ ≤\ (q\ :\ ℝ)\ *\ ‖x\ -\ y‖)}
\leanline{4}{(hfinite0\ :\ MemLp\ (fun\ ω\ ↦\ initial\ ω\ -\ Y\ 0\ ω)\ 2\ P)}
\leanline{4}{(hVfinite\ :\ ∀\ n,\ MemLp\ (V\ n)\ 2\ P)}
\leanline{4}{(hdefect\ :\ ∀\ n,\ (eLpNorm\ (V\ n)\ 2\ P).toReal\ ≤\ b)\ (n\ :\ ℕ)\ :}
\leanline{4}{wasserstein2Real}
\leanline{6}{(P.map\ (fun\ ω\ ↦\ ulaChain\ U\ h\ (initial\ ω)\ (fun\ k\ ↦\ noise\ k\ ω)\ n))\ π\ ≤}
\leanline{6}{(q\ :\ ℝ)\ \textasciicircum{}\ n\ *\ (eLpNorm\ (fun\ ω\ ↦\ initial\ ω\ -\ Y\ 0\ ω)\ 2\ P).toReal\ +}
\leanline{8}{b\ /\ (1\ -\ (q\ :\ ℝ))}
\end{leancode}
\leanlabel{How the Lean proof follows the argument}
\texttt{ula\_\allowbreak{}wasserstein\_\allowbreak{}bound\_\allowbreak{}of\_\allowbreak{}reference}: the intermediate facts include \texttt{X}, \texttt{e}, \texttt{hfinite}, \texttt{hc}, \texttt{hw}, in source order; the source steps apply or rewrite \texttt{hm.\allowbreak{}aestrongly\allowbreak{}Measurable}, \texttt{synchronous\_\allowbreak{}l2\_\allowbreak{}bound\_\allowbreak{}real}, \texttt{hw.\allowbreak{}trans}, \texttt{affine\_\allowbreak{}recurrence\_\allowbreak{}bound}, \texttt{q.\allowbreak{}coe\_\allowbreak{}nonneg}.
\end{leanbox}

\end{document}
